\section*{الفصل \ref{ch:g7:oxygen-in-the-water} --- الأكسجين في الماء}

\begin{solution}{exo:g7:oxygen-in-the-water:1}
نحو \qty{10}{mg} في اللتر --- أي أقل بنحو ثلاثين مرة مما يحمله لتر
من الهواء.
\end{solution}

\begin{solution}{exo:g7:oxygen-in-the-water:2}
الهواء فوقه، يذوب عند السطح --- وأسرع ما يكون حيث يُحرَّك الماء
ويُرشّ ويُرغّى؛ \omterm{def:g5:food-webs:roles}{والمنتِجون} تحته، يطلقون \omterm{def:g7:what-is-respiration:respiration}{الأكسجين} نهارًا --- وأسرع ما
يكون في ضوء ساطع مع نمو أخضر وفير.
\end{solution}

\begin{solution}{exo:g7:oxygen-in-the-water:3}
\omterm{def:g7:what-is-respiration:respiration}{تنفس} \omterm{def:g1:animals-around-us:animal}{الحيوانات}، \omterm{def:g7:what-is-respiration:respiration}{وتنفس} \omterm{def:g1:plants-around-us:plant}{النباتات} نفسها ليلًا، \omterm{def:g7:what-is-respiration:respiration}{وتنفس} المحلِّلين ---
وهو الذي يعظم حين تكثر المادة الميتة.
\end{solution}

\begin{solution}{exo:g7:oxygen-in-the-water:4}
الماء الدافئ يحمل أكسجينًا مذابًا أقل (فينخفض السقف) بينما يجري
محرّك كل متنفس أسرع (فيرتفع الطلب): عرض أقل وحاجة أكثر، معًا.
\end{solution}

\begin{solution}{exo:g7:oxygen-in-the-water:5}
نهارًا تضيف صناعة \omterm{def:g5:food-webs:roles}{المنتِجين} \omterm{def:g7:what-is-respiration:respiration}{الأكسجين} أسرع مما يستنزفه \omterm{def:g7:what-is-respiration:respiration}{التنفس} ---
فيصعد المنحنى إلى ذروة الأصيل. وطول الليل يسحب كل متنفس من المخزون
بلا صناعة البتة --- فينزلق المنحنى إلى أدناه قبيل الشروق مباشرة.
\end{solution}

\begin{solution}{exo:g7:oxygen-in-the-water:6}
الغني: السلمون المرقط (أو صغاره) \omterm{ex:g2:animals-grow-up:butterfly}{ويرقات} ذباب مايو تحت الحجارة.
والفقير: \omterm{def:g5:biodiversity:species}{الأنواع} المتحمّلة وحدها --- الشبّوط، وحلزونات البركة،
\omterm{ex:g2:animals-grow-up:butterfly}{واليرقات} ذوات الأنابيب التي تجلب الهواء من السطح.
\end{solution}

\begin{solution}{exo:g7:oxygen-in-the-water:7}
\omterm{def:g7:what-is-respiration:respiration}{تنفس} الليل --- \omterm{prop:g3:sorting-living-things:groups}{الأسماك} \omterm{def:g1:plants-around-us:plant}{والنباتات} ومحلِّلو الطين معًا --- يستنزف
البركة إلى دُنياها قبيل الشروق مباشرة، ولذلك تختنق أسماك البرك
المحمَّلة عند الفجر. والعلاجات: من المورّد الأول، حرّك السطح ---
شغّل نافورة أو شلالًا طول الليل؛ ومن المورّد الثاني، خفّف التشابك
والطين ليصغر \omterm{def:g7:what-is-respiration:respiration}{تنفس} الليل (وليبلغ إنتاج النهار الماءَ لا حصيرة السطح
وحدها).
\end{solution}

\begin{solution}{exo:g7:oxygen-in-the-water:8}
نفايات تدخل --- فيولم \omterm{def:g6:decomposers-and-soil:decomposers}{المحلِّلون} ويتكاثرون --- فيستنزف تنفسهم
المحتشد \omterm{def:g7:what-is-respiration:respiration}{الأكسجين} --- فتختفي \omterm{def:g5:biodiversity:species}{الأنواع} الطلوب ويستولي المتحمّلون. وفي
الأسفل تُستهلك النفايات أخيرًا، وتعيد الجداول الحصوية تهوية الماء،
فتعود \omterm{def:g5:biodiversity:species}{الأنواع} الطلوب: أُغلق المصرف، ولحق المورّدان.
\end{solution}

\begin{solution}{exo:g7:oxygen-in-the-water:9}
المِسبار يسجّل لحظة واحدة؛ أما السكان فيسجّلون تاريخًا. \omterm{def:g5:biodiversity:species}{فالنوع}
الطلوب لا يعيش في محطة إلا إذا بقي الماء جيدًا طول إقامته --- بما
في ذلك الفجرات والأسابيع السيئة التي لا تلتقطها زيارة واحدة.
والإحصاء دُنيا موسم كامل، مكتوبةً \omterm{def:g1:animals-around-us:animal}{بالحيوانات}.
\end{solution}

\begin{solution}{exo:g7:oxygen-in-the-water:10}
الخضرة كحساء البازلاء جمهرة هائلة من \omterm{def:g5:first-look-at-cells:cell}{الخلايا} المفردة الخضراء:
فهي تنتج نهارًا، لكن الحساء كله يتنفس كل ليلة دافئة --- والأسبوع
الحارّ الساكن يخفض السقف ويوقف إعادة الملء من السطح. فتتجه دُنيا
الفجر نحو الصفر: وهو طقس موت \omterm{prop:g3:sorting-living-things:groups}{الأسماك}. فشغّل النافورة أشدّ تشغيل
طول الليل والساعات الصغيرة --- أما النهار فيدبّر نفسه.
\end{solution}

\begin{solution}{exo:g7:oxygen-in-the-water:11}
السلمون المرقط من \omterm{def:g5:biodiversity:species}{الأنواع} الطلوب: فالأنهار الباردة تحمل السقف
العالي الذي يحتاج إليه؛ والسدود الصغيرة والمجاري الرشّاشة تُبقي
إعادة الملء دائمة في مقابل مصرف \omterm{prop:g3:sorting-living-things:groups}{الأسماك} المزدحمة. أما البحيرة
الدافئة الساكنة فتخفق في السقف وفي إعادة الملء معًا --- إنها بلاد
الشبّوط لا بلاد السلمون المرقط.
\end{solution}

\begin{solution}{exo:g7:oxygen-in-the-water:12}
العملة: كل متنفس تحت الماء يسحب من المخزون المذاب الواحد --- ومن
يزدهر وأين مسعَّر \omterm{def:g7:what-is-respiration:respiration}{بالأكسجين}
(\cref{prop:g7:oxygen-in-the-water:map}). وعرض النقد: الحرارة تضبط
كم يستطيع الماء أن يحمل --- فالبرد يرفع السقف والدفء يخفضه
(\cref{prop:g7:oxygen-in-the-water:drains}). وأكبر المنفقين: إذا
وُجدت مادة ميتة، فاق \omterm{def:g7:what-is-respiration:respiration}{تنفس} المحلِّلين المحتشد استنزافَ الجميع ---
وهو درس أنبوب معمل الألبان
(\cref{ex:g7:oxygen-in-the-water:waste}).
\end{solution}

\begin{solution}{pb:g7:oxygen-in-the-water:1}
\textbf{1.} م1 باردة --- سقف عالٍ --- ومخضوضة بالسدّ الصغير ---
فإعادة الملء من السطح بأقصى سرعة: فقراءات قرب السقف. وم2 دافئة ---
سقف منخفض --- وغير مخضوضة: فقراءة متواضعة كل ما يستطيع مورّداها
حمله.

\textbf{2.} \omterm{def:g5:biodiversity:species}{أنواع} م1 الطلوب تشهد \omterm{def:g7:what-is-respiration:respiration}{بأكسجين} غني ثابت؛ وطاقم م2
المتحمّل --- شبّوط، وحلزونات، وأصحاب أنابيب --- يعلّم ماءً يفتقر،
في ساعاته السيئة على الأقل.

\textbf{3.} \omterm{def:g7:what-is-respiration:respiration}{تنفس} المحلِّلين: فتصريف معمل الألبان الغني بالعضوي
يُطعم جمهرة محتشدة منهم، ووليمتهم تستنزف المخزون إلى ما دون ما
تفسّره الحرارة وحدها بكثير.

\textbf{4.} كيلومترا استهلاك استنفدا النفايات --- فانتهت الوليمة
وخفّ المنفقون --- بينما أعاد الجدول الحصوي خضّ \omterm{def:g7:what-is-respiration:respiration}{الأكسجين} إلى
الداخل: فاستُعيد المورّدان، وأُغلق المصرف، والإحصاء يشهد بالتعافي.

\textbf{5.} موجة يومية: قاع عند الفجر عند \qty{4}{mg/L}، وذروة في
الأصيل عند \qty{9}{mg/L}. والمتبادلان: صناعة النهار (تضيف نهارًا)
\omterm{def:g7:what-is-respiration:respiration}{والتنفس} على مدار الساعة (يستنزف وحده طول الليل).

\textbf{6.} برد الشتاء يرفع السقف فيحمل الماء أكثر؛ والبرد يبطّئ
كل متنفس --- \omterm{prop:g3:sorting-living-things:groups}{الأسماك}، \omterm{def:g1:plants-around-us:plant}{والنباتات}، ومحلِّلي الطين --- فتصغر المصارف.
وأثبت أيضًا، لأن موجة \omterm{def:g5:food-webs:roles}{المنتِجين} اليومية أصغر.

\textbf{7.} م3 أولًا (مستنزفة منذ الآن، والحرّ يسرّع محلِّليها)؛
وم2 ثانيًا (سقف منخفض، وماء ساكن، وقيعان عند الفجر)؛ وم4 ثالثًا
(تتعافى لكنها أسفل المشكلة)؛ وم1 أخيرًا (باردة مخضوضة --- فالسدّ
الصغير يواصل إعادة الملء مهما يكن الطقس).

\textbf{8.} ما قبل الشروق هو الدُّنيا اليومية: فالقراءات عندئذ
تسجّل أسوأ ما على حياة النهر أن تنجو منه. أما جدول الأصيل فسيحفظ
السنة تحت أحسن لحظاتها.

\textbf{9.} المِسبار يتعافى في أسابيع --- فالمصرف يُغلق حين يتوقف
غذاء المحلِّلين. والإحصاء يتعافى أبطأ وبترتيب: فيخفّ الطاقم
المتحمّل أولًا، ثم يعيد المستوطنون المنجرفون والطائرون
(بمنطق \cref{ch:g6:how-plants-spread}، في نسخته الحيوانية) بذر
\omterm{def:g5:biodiversity:species}{الأنواع} الطلوب --- \omterm{ex:g2:animals-grow-up:butterfly}{يرقات} ذباب مايو أولًا من الأعلى، والسلمون
المرقط أخيرًا.

\textbf{10.} الظل يُبقي المقطع باردًا أيضًا --- والبرد هو السقف.
فالقطع سيضيف إنتاج النهار لكنه سيخفض السقف ويسرّع كل مصرف في الماء
الأدفأ؛ وفي مقطع بطيء يرجّح أن تفوق الخسارة الكسب. والجواب: أبقِ
الأشجار، وحرّك الماء بدل ذلك.

\textbf{11.} الإحصاء: \omterm{ex:g2:animals-grow-up:butterfly}{فيرقات} ذباب مايو لا تستطيع أن تزوّر موسمًا
جيدًا --- فقد عاشت كل فجر منه. أما قراءة الأصيل المتوسطة فلحظة
واحدة، ولعلها يوم سيئ؛ \omterm{ex:g2:animals-grow-up:butterfly}{واليرقات} سجلّ الموسم المشفوع باليمين.

\textbf{12.} ``أبقِ النهر باردًا مخضوضًا غير محمَّل، فيدبّر
\omterm{def:g7:what-is-respiration:respiration}{الأكسجين} --- وكل ما يتنفسه --- نفسه بنفسه.''
\end{solution}
